\scanpage{62}
% archive sheet number: 61

\unsure{Celà|Ceci|Cela} \unsure{marche|mène} aussi si \deleted{\unsure{$Z$|$X$}}\inserted{$(\unsure{q_2|q_{0}|\psi_1})$} $Y$ \inserted{cpct} admet \deleted{des} rec.ts arbitrairement fins \inserted{finis} par des $Y_i$ \inserted{fermés ou ouverts} tels que les intersections $Y_{i_1}\cap Y_{i_2}\cap\ \cap Y_{i_{\unsure{r|\nu}}}$ $(\unsure{r|\nu}\geq1)$ de tels $Y_i$ soient $\in\unsure{\mathcal E|\mathscr E}_{\chi!}$ \ldots.

Dans le cas \unsure{du|des} $\chi$, on utilise la suite spectrale d'\uline{homologie}, on est ramené :

\uline{Lemme} Soit $R$ syst. \unsure{loc$^{\mathrm{t}}$|local} de \unsure{$k$-vectoriels|$k$-vectoriel|homotopique} sur $Y\in\mathcal E_\chi$, \unsure{fibres|fibrés} de rang \inserted{fini} $b$. Alors

\[
\chi(Y,R)\overset{\mathrm{d\acute ef}}{=}\sum_p(-1)^p
\underset{\text{le dual est }H^p(Y,\check R)}{\underline{H_p^{\deleted{\unsure{p|\flat}}}(Y,R)}}
\]

existe et

\[
\chi(Y,R)=b\chi(Y)
\]

Cela montre pourvu qu'on fasse l'hyp. \deleted{\unsure{(q$_2$)}} \inserted{que $Y$} est $\sim Y'$ satisfaisant la condition $(\unsure{q_2|q_{0}|\psi_2})$
\begin{center}$=$\end{center}

\deleted{On va se borner}

\textcircled{\unsure{4|h}} Définissons $\mathcal C_{-1}={}$ ens. réduit : l'espace vide

et si $n\geq0$ $\mathcal C_n^{\deleted{\unsure{\mathrm{cpct}}}}={}$ \deleted{catégorie} \inserted{ens.} des espaces cpcts \inserted{$X$} tels que tt pt \inserted{$x\in X$} ait un voisinage $C_x$ tel qu'il existe un espace compact $S\in\mathcal E_{n-1}$ et un homéom. $\mathrm{Cône}(S)\simeq C_x$ transformant le sommet du cône en $x$.

\margin{\uline{NB} $\mathcal E_0$ espaces discrets\\
$\mathcal E_1$ 1-complexes topologiques\\
$\mathcal E_2$ espaces loc. triangulables de dim $\leq2$ [ils sont \unsure{tri-angul.|triangulables} \unsure{ds|dans} le cas cpct]\\
$\mathcal E_3$ espaces loc. triangulables de dim $\leq3$ [\unsure{ils sont tr. triangul.|triangulables} \unsure{ds|dans} le cas cpct]\\
$\mathcal E_4$ esp. loc. \unsure{\textquotedblleft id\textquotedblright} dim $\leq4$ (Moïse) S'ils sont cpcts, il n'est pas clair qu'ils sont triangulables \ldots.}

\[
\mathcal C_\infty=\bigcup\mathcal C_n\qquad\text{(espaces loc. côniques)}
\]

$\mathcal E_\infty={}$ \deleted{\unsure{esp}} des espaces loc. cpcts \inserted{séparés} $X$ tels que $\hat X$ soit loc. cônique i.e. soit $\in C_\infty$.

\[
\text{\uline{Prop}}\qquad \mathcal C_\infty\subset\mathcal E_\chi\cap\mathcal E_{\chi!}
\]

Soit $X\in\mathcal E_\infty$. Si $X$ est cpct, \inserted{donc $X\in C_\infty$} comme $X$ est loc. \unsure{contr. et il|contr. et $L$|contractile}, on sait que $X\in\mathcal E_\chi$, donc $X\in\mathcal E_{\chi!}$. \deleted{\unsure{Cas $X$ quelconque}} \inserted{Sinon} : $X$ loc. cpct $=\hat X-(\omega)$, avec $\hat X\in\mathcal C_\infty$, on a $\hat X\in\mathcal E_\chi$, donc $\hat X-\omega=X\in\mathcal E_{\chi!}$ par (b). Pour prouver $X\in\mathcal E_\chi$, prenons un $C_x\simeq C(S)$, prenons un homothétique $\frac12 C_x=T$, et
